\section{Implementation}	\label{a:ab_rts_en}

This chapter covers the implementation of Bayesian optimization with a Cheetah based physics informed prior for the transverse beam tuning at the ARES Experimental Area. 
Instead of being one large script, the implementation code is systematically divided into python modules, each having their distinct functions:
\begin{itemize}
    \item \textbf{\texttt{bo\_cheetah\_prior\_ares.py}} includes the simulated ARES environment \texttt{ares\_problem()} and the GPyTorch compatible prior mean class \texttt{AresPriorMean} with its trainable  misalignment parameters. This is where most of the physics lives.
    \item \textbf{\texttt{eval\_ares.py}} constructs the optimizers (Standard BO, BO with cheetah prior and Nelder-Mead) along with the experimental scenarios (matched, mismatched and matched\_prior\_newtask) and loops through the optimization trials. 
    \item \textbf{\texttt{eval\_ares\_metrics.py}} in accordance with the evaluation process \cite{kaiser2024rl}, it calculates the performance metrics used to compare the techniques. 
    \item \textbf{\texttt{plot\_ares\_results.py}} generates plots for the results. 
\end{itemize}

The plotting and metrics modules are then connected by the fifth script \textbf{\texttt{run\_evaluation.py}} to run altogether, it can also generate result tables in LaTex format. Table \ref{table:libraries} shows the summary of primary libraries used. 

\begin{table} [htbp]
    \centering
    \small
    \begin{tabular}{p{2.5cm} p{5.5cm} p{1.5cm}}
        \toprule
        \textbf{Libraries} & \textbf{Purpose} & \textbf{Version} \\
        \midrule
        \texttt{Cheetah} & PyTorch based differentiable beam tracking simulator. & 0.8.0 \\
        \texttt{GPyTorch} & Gaussian process framework, base class for prior mean module & 1.14.2 \\ 
        \texttt{Xopt} & BO loop, UCB acquisition Nelder-Mead baseline. & 2.6.7 \\ 
        \texttt{BoTorch} & BO primitives used internally by Xopt & 0.16.0 \\ 
        \bottomrule
    \end{tabular}
    \caption{Primary libraries/frameworks used.}
    \label{table:libraries}
\end{table}

\subsection{Initial Setup}
Before running any optimizers, we need something to optimize against. Defined in \texttt{bo\_cheetah\_prior\_ares.py}, 
the method \texttt{ares\_problem()} fulfills this purpose. A beam is propagated across a Cheetah model of the ARES experimental 
area using a set of magnet parameters, and the beam quality is then returned. In practice, every call to this function 
mimics what would happen if the user adjusted the magnet settings on the real machine and then looked at the diagnostic screen. \

In this project, cheetah's ParameterBeam is used to describe the beam entering the ARES Experimental Area. \texttt{ParameterBeam} was used 
simply because it is faster, which enables it to run the experiments more quickly and iterate on the implementation with shorter turnaround times.
 Also note that Cheetah's \texttt{ParticleBeam} could have also been just as well but with more computational cost. Since speed is not the most critical 
 concern in this project, there is nothing holding us back from using either representation for this application. The ARES experimental area only 
 involves linear optics, so both approaches get equivalent results. These values match the typical ARES operating conditions, with a beam energy of 
 around 107 MeV and transverse sizes of 200 $\mu$m.

A JSON file describing the full ARES accelerator is used to load the lattice itself. The code extracts the sub-segment between AREASOLA1 (marker) and AREABSCR1 (screen) 
as we only require the experimental ares. Listing \ref{code:ares_loading} demonstrates the loading of the lattice and the assignment of the five magnet parameters prior to tracking.


\lstset{language=Python,
    basicstyle=\ttfamily\small,      % typewriter font, smaller size
    keywordstyle=\color{blue},        % keywords in blue
    commentstyle=\color{teal},        % comments in green
    stringstyle=\color{purple},       % strings in purple
    showstringspaces=false,           % don't mark spaces in strings
    breaklines=true,                  % automatic line breaking
    numbers=left,                      % line numbers on left
    numberstyle=\tiny\color{gray},     % line number style
    frame=single,                       % add a frame around the code
    backgroundcolor=\color{lightgray!20}, % very light gray background
    xleftmargin=2em,                    % left margin for line numbers
    framexleftmargin=2em                 % extend frame to include numbers}
}

\begin{lstlisting}[caption={ARES lattice loading and magnet parameters configuration.}, label={code:ares_loading}]

# Load ARES lattice
ares_segment = cheetah.Segment.from_lattice_json("ARESlatticeStage3v1_9.json")
ares_ea = ares_segment.subcell("AREASOLA1", "AREABSCR1")

# Set magnet parameters
ares_ea.AREAMQZM1.k1 = torch.tensor(input_param["q1"])
ares_ea.AREAMQZM2.k1 = torch.tensor(input_param["q2"])
ares_ea.AREAMCVM1.angle = torch.tensor(input_param["cv"])
ares_ea.AREAMQZM3.k1 = torch.tensor(input_param["q3"])
ares_ea.AREAMCHM1.angle = torch.tensor(input_param["ch"])

if misalignment_config is not None: 
    for magnet_name , (dx, dy) in misalignment_config.items():
        magnet = getattr(ares_ea, magnet_name)
        magnet.misalignment = torch.tensor([dx, dy], dtype=torch.float32
        )
else:
        ares_ea.AREAMQZM1.misalignment = torch.tensor([0.0, 0.0])
        ares_ea.AREAMQZM2.misalignment = torch.tensor([0.0, 0.0])
        ares_ea.AREAMQZM3.misalignment = torch.tensor([0.0, 0.0])

out_beam = ares_ea(incoming_beam)
\end{lstlisting}

For the three quadrupole strengths ($k^{Q1}$, $k^{Q2}$, $k^{Q3}$) and corrector angles ($\theta^{CV}$ and $\theta^{CH}$),  we have limited the range from physical magnet's 
$\pm72 m^{-2}$ to our simulation's $\pm30 m^{-2}$, and from $\pm 6.2 mrad$ to $\pm 6.0 mrad$ respectively because it does not make sense since there is no reasonable 
way for them to go that far out. The code iterates through the dictionary and applies each ($\Delta_x, \Delta_y$) pair to the corresponding 
quadrupoles when a misaligned configuration is provided. This enables us to model the realistic scenario in which magnets are not perfectly 
centred on the beam. To let the optimizer know how good or bad the result is, we need a single number. Following the previous ARES studies approach \cite{kaiser2024bridging}, we use the mean absolute error for each of the four beams observable.

\begin{equation}
\label{mae_repeated}
 mae = 0.25(|\mu_x|+\sigma_x+|\mu_y| + \sigma_y)   
\end{equation}


Here in equation \ref{mae_repeated}, $\mu_x$ and $\mu_y$ are the beam positions (ideally, zero means beam is centred) and $\sigma_x$ and $\sigma_y$ are beam sizes (the smaller the better). Combining all four into a single metric forces the optimiser to center and focus the beam simultaneously. 
The implementation is straight forward. The function also returns the individual beam parameters in addition to MAE so that they can be logged for later analysis.

\subsection{Physics Informed Prior Mean Module}
The heart of this project work is the \texttt{AresPriorMean}, defined in \texttt{bo\_cheetah\_prior\_ares.py}. 
Its role is to turn the cheetah simulation into a format that GPyTorch can use as the mean function of the Gaussian Process. 
Instead of using GP prior (typically with zero mean assumption), it begins with the best estimates from the physics model and 
just needs to learn the residual which is the difference between what the simulator predicts and what is actually observed. \

\texttt{AresPriorMean} class specifically inherits from \texttt{gpytorch.means.Mean} and implements a \texttt{forward()} method. 
Everytime GPyTorch needs the prior mean at a set of input points, it calls this method. The remaining GP components, such as kernel evaluation, 
posterior conditioning and acquisition function optimization, perform exactly like they would with any other mean function. 
Another important aspect is that Xopt builds the input tensor for the GP by sorting variables in alphabetical order, which means that since our 
magnets variables are named $q1$, $q2$, $cv$, $q3$ and $ch$, column 0 of the tensor would  always be $ch$, not $q1$ as one might expect from the beamline order. 
In order to avoid silently feeding the wrong values to incorrect magnets, the class defines explicit index constants (Listing \ref{code:xopt_indexing}).

\begin{lstlisting}[caption={Explicit indexing to ensure correct mapping in Xopt input tensor.}, label={code:xopt_indexing}]

# Load ARES lattice
class AresPriorMean(Mean):
    # Considering alphabetical order
    VAR_ORDER = ['ch', 'cv', 'q1', 'q2', 'q3']  
    IDX_CH = 0
    IDX_CV = 1
    IDX_Q1 = 2
    IDX_Q2 = 3
    IDX_Q3 = 4
\end{lstlisting}

As discussed earlier, in real accelerators, the quadrupole magnets are never perfectly aligned. A magnet that is even with a fraction of a millimeter, 
divertes the beam in a way that would not be predicted by a perfectly aligned simulation, partially acting as a steering factor.  
In order to handle this, the prior module treats the quadrupole misalignments as learnable parameters ($\Delta_{xQ1}$, $\Delta_{yQ1}$, $\Delta_{xQ2}$, $\Delta_{yQ2}$, $\Delta_{xQ3}$, $\Delta_{yQ3}$).
Each of the three quadrupoles has horizontal and vertical misalignment values, making a total of six values. In order to do online system identification, GPyTorch’s marginal likelihood maximization 
adjusts these values in addition to the standard GP hyperparameters as the optimizer collects data. Three coordinated layers are used in GPyTorch to register each misalignment parameter, these three 
layers are summarized in table \ref{three_layers}. On top of this, $@property$ accessors offer a convenient interface, reading a misalignment property applies the forward transform (raw to constrained) and when 
it is written, the inverse transform is applied. This three layer pattern is repeated identically for each of the six parameters.  The constructor is fairly lengthy because of this six times repeated 
pattern, although a loop based factory would have been cleaner, but we go with the explicit version since it has the advantage of being easy to read and to debug. 

\begin{table} [htbp]
    \centering
    \small
    \begin{tabular}{p{4.0cm} p{4.0cm} p{7.0cm}}
        \toprule
        \textbf{Component} & \textbf{GPyTorch API} & \textbf{What it does} \\
        \midrule
        Raw Parameter & \texttt{register\_parameter()} & A raw unconstrained \texttt{nn.Parameter} is created at first, which is what PyTorch’s optimiser updates directly.\\
        Interval Constraint & \texttt{register\_constraint()} & Secondly, an interval constraint is used that maps this raw value through a sigmoid based transformation into the physically meaningful range of $\pm0.5mm$, this serves as the bound guaranteeing that the learnt offset can never be greater than what is practical for quadrupole misalignments. \\ 
        Smoothed box prior & \texttt{register\_prior()} & Third, a \texttt{SmoothBoxPrior} is registered across the confined space, adding a soft regularization on top of the hard constraint. This helps the optimizer avoid getting stuck at the boundaries by gently penalizing values near the edges while assigning a fairly uniform probability throughout the interval. \\ 
        \bottomrule
    \end{tabular}
    \caption{Three layered pattern to register misalignments as trainable parameters.}
    \label{three_layers}
\end{table}

When GPyTorch calls \texttt{forward()}, the method recieves a tensor X whose last dimension contains 5 columns (the five magnets, in alphabetical order), the steps include:
\begin{enumerate}
    \item Using the index constants from Listing \ref{code:xopt_indexing}, slice out each magnet parameter. 
    \item Assign them to corresponding components of cheetah lattice.
    \item Apply the current misalignment values to the three quadrupoles by stacking them.
    \item Run the beam through the lattice
    \item Return the MAE (equation \ref{mae_repeated}).
\end{enumerate}
Since all of these operations (the beam tracking, the absolute values, and the summation) are native PyTorch operations, the full forward pass is differentiable. 
Gradients could flow backward through the simulator and into the six misalignment parameters, which is what allows the joint optimization of both misalignment 
values as well as usual GP hyperparameters.

\subsection{Optimization Runner and Evaluation Metrics}

After setting up the environment and the prior setup, the remaining part is the runner that actually executes the experiments. This is handled by an \texttt{eval\_ares.py} file 
that takes arguments for the experimental scenarios (matched, mismatched, matched\_prior\_newtask), the number of trials, the evaluation budget and the type of optimizer (BO, BO\_prior, NM) 
where BO is the standard Bayesian Optimization, BO\_prior is Bayesian Optimization with physics prior and NM is Nelder-Mead. Xopt defines the optimization problem declaratively using a 
YAML based specification known as VOCS (Variable and Objective Configuration Specification). Listing \ref{code:VOCS_configuration}  shows the five variables and the single objective. Where the unit for quadrupole magnet
ranges are $m^{-2}$, for corrector magnets ranges are $rad$ and for objective (mae) is metre($m$)


\begin{lstlisting}[caption={VOCS configuration.}, label={code:VOCS_configuration}]
vocs_config = """
        variables:
            q1: [-30, 30]
            q2: [-30, 30]
            cv: [-0.006, 0.006]
            q3: [-30, 30]
            ch: [-0.006, 0.006]
        objectives:
            mae: minimize
    """
    vocs = VOCS.from_yaml(vocs_config)
\end{lstlisting}

To test the prior under various conditions, the runner supports three scenarios given below, 
they are also summarized in Table \ref{table:scenarios_summary}:

\begin{enumerate}
    \item \textbf{Matched Scenario:} This is the most simple and ideal scenario where the beam and misalignments in the environment are exactly what the prior expects, the default values (including all misalignment values to be zero), thus the physics model perfectly matches reality. This scenario exists mainly to verify that the implementation works correctly and to see how much the prior helps when it is perfectly accurate since there is no model reality gap at all. In practice, this situation never happens on a real machine. 
    \item \textbf{Mismatched Scenario:} In this scenario, a gap is introduced on purpose to create mismatch between the prior and the reality, a different beam ($\sigma_y$ changed from 200$\mu$m to 100$\mu$m), and non zero misalignment values over all three quadrupoles are used in the environment(Table \ref{table:misalignment_groundtruth}). The prior, however, starts from the default beam and zero misalignments as it does not know about any of these changes. We tell Xopt to include the six misalignment parameters in the GP’s marginal likelihood optimisation by setting \texttt{trainable\_mean\_keys=["mae"]} in the StandardModelConstructor so they can be learned from data as it is received. This is the scenario that tests whether online systems identification actually works. 
     \begin{table} [htbp]
    \centering
    \small
    \begin{tabular}{p{4.0cm} p{2.0cm} p{2.0cm}}
        \toprule
        \textbf{Magnet} & \textbf{$\Delta_x$ (mm)} & \textbf{$\Delta_y$ (mm)} \\
        \midrule
        AREASMQZM1 & 0.000 & +0.200 \\
        AREASMQZM2 & +0.100 & -0.300 \\ 
        AREASMQZM3 & -0.100 & +0.150 \\ 
        \bottomrule
    \end{tabular}
    \caption{Ground truth quadrupole misalignments.}
    \label{table:misalignment_groundtruth}
\end{table}


    \item \textbf{Matched\_prior\_newtask Scenario:} The simulated accelerator setup is the same as in the mismatched scenario (modified beam and non-zero misalignments), but this time the prior is initialised with the correct values which are the right beam parameters and the true misalignment values. So this way, we are handing the prior a perfectly adjusted physics model and asking how fast can the optimizer converge when it already knows the true system’s state, compared to having to learn it during the process.
\end{enumerate}

For the evaluation in this project, the matched\_prior\_newtask serves as our actual matched case, since it uses the realistic 
accelerator setup having modified beam and misalignments rather than the ideal zero misalignment defaults. We compare four methods, 
Nelder-Mead, Standard BO (zero mean), BO with Cheetah prior (mismatched scenario) and BO with Cheetah prior (matched\_prior\_newtask scenario). 
All four approaches face the same simulated accelerator so that the comparison is fair,  what differs between them is only the optimization strategy, 
and in case of prior informed methods, how much the prior knowledge is available from the start for each method. 

   
\begin{table} [htbp]
    \centering
    \small
    \begin{tabular}{p{2.5cm} p{2.0cm} p{2.0cm} p{2.0cm} p{2.0cm}}
        \toprule
        \textbf{Scenario} & \textbf{Modified Beam} & \textbf{True Misalignments} & \textbf{Prior Misalignment} & \textbf{Trainable $\phi$} \\
        \midrule
        Matched & --- & --- & --- & --- \\
        Mismatched & $\checkmark$ & $\checkmark$ & 0 & $\checkmark$ \\ 
        Matched prior (new task) & $\checkmark$ & $\checkmark$ & = true & ---  \\ 
        \bottomrule
    \end{tabular}
    \caption{Three experimental scenario summary.}
    \label{table:scenarios_summary}
\end{table}

In order to have a fair comparison, every trial starts from the same initial points:

\begin{equation}
    x_{0} = (k^{Q1}=10.0, k^{Q2}=-10.0, \theta^{CV} = 0.0, k^{Q3} = 10.0, \theta^{CH} = 0.0)
\end{equation}

After this initial evaluation, the runner loops through the \texttt{xopt.step()} for a fixed number of iterations. 
The default in the code is 50 steps, but for the evaluation presented in this report, we used 100 steps to give each 
method more room to converge and to better observe the later stage behaviors of the optimizers. Each step involves 
fitting the GP to all the data collected so far, selecting the next point by optimizing the UCB acquisition function 
and testing that point on the simulated accelerator. In the case of Nelder-Mead, the same \texttt{step()} interface 
updates the simplex instead. For mismatched scenario when the prior’s misalignment parameters are trainable, fitting 
the GP also involves optimizing over \textbf{$\phi$}. As a result, each new data point gives the marginal likelihood 
optimiser more information about where the magnets are actually positioned and the prior gradually gets more accurate 
which is the online system identification already discussed above. After each trial completes, the script stores the 
full MAE history, the overall best, and the learned misalignment values (where applicable). With a fixed random seed, 
ten separate trials are conducted per configuration and the results are saved to CSV files for further analysis. 




\subsection{Evaluation Metrics}
The evaluation code is located in \texttt{eval\_ares\_metrics.py} that follows to \cite{kaiser2024rl}. 
The metrics computed for every trial are as follows:
\begin{enumerate}
    \item \textbf{Final error:} is the best (minimum) MAE reached at any point during the trial.
    \item \textbf{RMSE:} is the root mean square of the overall best MAE curve that captures how quickly the optimizer improves and not just the final value.
    \item \textbf{Improvement percentage(\%):} is the improvement measured from initial to best $(((\text{MAE}_0 - \text{MAE}_{\text{best}})/\text{MAE}_0) \times 100)$.
    \item \textbf{Steps to target:} shows how many evaluations does it takes to drop below 40 $\mu$m of MAE (which is the threshold set) for the first time.
    \item \textbf{Success rate:} is the percentage of trials that reach the 40 $\mu$m target within the budget.
    \item \textbf{Steps to convergence:} shows the first step after which the overall best varies by less than $10^{-6}$ .
\end{enumerate}

Though all the metrics provide quite good insights, the final error (the best MAE achieved) is the most crucial one for the practical  purposes, 
and the steps to target are also informative because it tells how quickly the optimiser gets the beam into an acceptable state. \texttt{plot\_ares\_results.py} 
produces plots and tables, styled to match in \cite{kaiser2024bridging} and \cite{kaiser2024rl}, which can then be seen later under the "Result and Discussion" sections.  